%%%%%%%%%%%%%%%%%%%%%%%%%%%%%%%%%%%%%%%%%%%%%%%%%%%%%%%%%%%%%%%%%%%%%%%%%%%%%%%%
\section{Snakemake Nanopub Reports}
{   
    \usebackgroundtemplate{
        \vbox to \paperheight{\vfil\hbox to \paperwidth{\hfil\includegraphics[width=\paperwidth]{reports/nanopub_workflow_report.png}\hfil}\vfil}
    }
    \frame{
        \frametitle{\Snakemake Reports with Nanopubs}
        \begin{mdframed}[tikzsetting={draw=white,fill=white,fill opacity=0.8,
                line width=0pt},backgroundcolor=none,leftmargin=0,
            rightmargin=150,innertopmargin=4pt,roundcorner=10pt]
            \tableofcontents[currentsection,sections={1-4},hideothersubsections]
        \end{mdframed}
        \vfill\vspace{16mm}\hfil \hfill{\tiny Screenshot -- upper section of a registered workflow nanopub}
    }
}

%%%%%%%%%%%%%%%%%%%%%%%%%%%%%%%%%%%%%%%%%%%%%%%%%%%%%%%%%%%%%%%%%%%%%%%%%%%%%%%%
\begin{frame}
    \frametitle{\HandsOn{How to register your Workflow}}
    We created a \lhref{https://w3id.org/np/RAOT7z3RA0XYlHIikne8rfUUYZrtHyrzXBD1HpI_GvcRk}{registration template}: \altverb{RAOT7z3RA0XYlHIikne8rfUUYZrtHyrzXBD1HpI_GvcRk}
    \begin{task}
       Use this template to register a workflow. Possibly \emph{your} workflow? Else the tutorial workflow.\\
       To register a workflow, it needs to be registered with a DOI (or a nanopub itself):
       \begin{itemize}
           \item On Zenodo you can create a DOI via GitHub, e.g. \lhref{https://doi.org/10.5281/zenodo.17087101}{this one}.
           \item Alternatively, \lhref{https://workflowhub.eu/search?q=Snakemake\#workflows}{WorkflowHub} has a few registered workflows -- we will work to extend this.
           \item Or you use \lhref{https://doi.org/10.14279/eceasst.v83.2600}{the DOI for a tutorial workflow}  registered with the HPC teaching alliance for \Snakemake.
       \end{itemize}
    \end{task}
\end{frame}

%%%%%%%%%%%%%%%%%%%%%%%%%%%%%%%%%%%%%%%%%%%%%%%%%%%%%%%%%%%%%%%%%%%%%%%%%%%%%%%%
\begin{frame}
    \frametitle{What if a workflow is registered multiple times?}
    \begin{docs}
        The Nanopublications framework does not prevent this. It is not desirable, but cannot be avoided.
        \begin{itemize}[<+->]
            \item You can retract accidentally created nanopubs (or those with a glitch).
            \item Search before registering \lhref{https://nanodash.knowledgepixels.com/query?id=RApMWbZ0ixGj88tPB4cEamKJQuQTZ5Dxpjh1a5BWteTdE/SnakemakeWorkflowRegistrations}{using this query}. 
        \end{itemize}
    \end{docs}
    \pause
    \begin{hint}
        We will retract unnecessarily created nanopubs, e.g. from the example DOIs.
    \end{hint}
\end{frame}

%%%%%%%%%%%%%%%%%%%%%%%%%%%%%%%%%%%%%%%%%%%%%%%%%%%%%%%%%%%%%%%%%%%%%%%%%%%%%%%%
\begin{frame}[fragile]
    \frametitle{\HandsOn{Using the plugin to create your Nanopub report}}
    First we need to install it into a working \Snakemake environment:
    \begin{lstlisting}[language=Bash, style=Shell]
$ # pip or mamba
$ pip install snakemake-report-plugin-nanopub
    \end{lstlisting}
    Then take your just created nanopub to register your workflow:
    \begin{lstlisting}[language=Bash, style=Shell]
$ snakemake ... --reporter nanopub \
> --report-nanopub-workflow-id <nanopub workflow id>
    \end{lstlisting}
    \begin{warning}
        This may trigger a workflow to run!
    \end{warning}
\end{frame}

%%%%%%%%%%%%%%%%%%%%%%%%%%%%%%%%%%%%%%%%%%%%%%%%%%%%%%%%%%%%%%%%%%%%%%%%%%%%%%%%
\begin{frame}[fragile]
    \frametitle{\HandsOn{Browsing the report}}
    Browse your own nanopub report or \lhref{https://w3id.org/np/RAK9xz_ccnu0Xhs4vX2KtqCxX44mmSt6nq-ePLeewMrFE}{this one}
    \begin{hint}
        Basically you can use this to simplify a materials \& methods section of a paper like this:\\
        
        
    \end{hint}
\end{frame}

%%%%%%%%%%%%%%%%%%%%%%%%%%%%%%%%%%%%%%%%%%%%%%%%%%%%%%%%%%%%%%%%%%%%%%%%%%%%%%%%
\begin{frame}
    \frametitle{An example knowledge graph}
    To visualize the result we can use the plugin to create a graph like this:
    \begin{centering}
    \includegraphics[width=0.9\textwidth]{reports/example_knowledgegraph.png}
    \end{centering}
    \begin{hint}
        It will contain the descriptions the workflow designer (you?) entered into the workflow, the dataset, etc..
    \end{hint}
\end{frame}

%%%%%%%%%%%%%%%%%%%%%%%%%%%%%%%%%%%%%%%%%%%%%%%%%%%%%%%%%%%%%%%%%%%%%%%%%%%%%%%%
\begin{frame}[fragile]
    \frametitle{Creating this knowledge graph}
    The nanopub reporter plugin for \Snakemake has a script on board:
    \begin{lstlisting}[language=Bash, style=Shell]
$ plot-nanopub-knowledge-graph \
> --dataset-nanopub-id <dataset_id> \
> --workflow-nanopub-id <id of registered workflow> \
> --workflow-configuration-id <id of registered conf> \
> --report-nanopub-id <id of registered report> \ 
> -o example_knowledgegraph.png
    \end{lstlisting}
    \begin{hint}
        \begin{itemize}
            \item file type determined by suffix (\altverb{png}, \altverb{pdf}, \altverb{svg})
            \item all flags are mandatory
        \end{itemize}
    \end{hint}
\end{frame}
